% GENERATED FILE. Edit canonical lessons, then run npm run generate:books.
% canonical-source-hash: fnv1a64:ba03ede7f7779ee2
% canonical-lessons: RU-C490-vyzdoravlivat, RU-C490-vyzdorovleniye, RU-C490-popravlyatsya, RU-C490-glotat, RU-C490-sirop

\chapter{Getting Better}
\label{ch:ru-getting-better}
\hlchaptermodality{490}

\begin{chapteropening}
\textbf{By the end of this chapter:} \emph{I can talk about recovering after an illness.}

Complete the last lesson of chapter 490: I can talk about recovering after an illness.
\end{chapteropening}

\section[\texorpdfstring{\ru{выздоравливать} (vyzdorávlivat)}{выздоравливать (vyzdorávlivat)}]{\ru{выздоравливать} (vyzdorávlivat) — to recover}

\label{lesson:RU-C490-vyzdoravlivat}



\begin{quote}
Before the new one: say the Russian for an X-ray, then the Russian for a surgeon.
\end{quote}

\subsection*{You'll want to know: \ru{выздоравливать}}
\textbf{\ru{выздоравливать}} (\emph{vyzdorávlivat}) — \textquotedblleft{}to recover\textquotedblright{}.

\begin{quote}
\ru{Мой} \ru{брат} \ru{медленно} \ru{выздоравливает} \ru{после} \ru{гриппа}.
\end{quote}

\begin{quote}
\emph{Moy brat médlenno vyzdorávlivayet pósle gríppa.}
\end{quote}

\begin{quote}
My brother is slowly recovering from the flu.
\end{quote}

Imperfective; perfective \ru{выздороветь} (\ru{выздоровею}).

\begin{cousinweb}
\emph{\ru{выздоравливать}} is built on \emph{\ru{вы}-} and \emph{\ru{здоровый}} 'healthy'; EtymDB lists \emph{\ru{выздороветь}} and \emph{\ru{здоровый}} without a source.
\end{cousinweb}

\begin{grammarlens}[title={What you've built}]
One more word for telling, explaining and describing at B1.
\end{grammarlens}

\subsection*{Your turn}
\begin{itemize}
\raggedright
  \item \emph{Say it:} \emph{vyzdorávlivat}
  \item \emph{Say it:} \emph{vyzdorávlivat}, once more
  \item \emph{Say it:} a sentence of your own with \emph{vyzdorávlivat}
  \item \emph{Recall:} say the Russian for an X-ray, then the Russian for a surgeon, then say \emph{vyzdorávlivat} again
\end{itemize}

\begin{tcolorbox}[breakable,colback=teal!4,colframe=teal!35!black,title={Before you move on}]
What does \ru{выздоравливать} mean? (\textquotedblleft{}To recover\textquotedblright{}.) Say it once more.
\end{tcolorbox}
\section[\texorpdfstring{\ru{выздоровление} (vyzdorovléniye)}{выздоровление (vyzdorovléniye)}]{\ru{выздоровление} (vyzdorovléniye) — recovery}

\label{lesson:RU-C490-vyzdorovleniye}



\begin{quote}
Before the new one: say the Russian for a surgeon, then the Russian for to recover.
\end{quote}

\subsection*{You'll want to know: \ru{выздоровление}}
\textbf{\ru{выздоровление}} (\emph{vyzdorovléniye}) — \textquotedblleft{}recovery\textquotedblright{}.

\begin{quote}
\ru{Желаем} \ru{вам} \ru{скорого} \ru{выздоровления}!
\end{quote}

\begin{quote}
\emph{Zheláyem vam skórogo vyzdorovléniya!}
\end{quote}

\begin{quote}
We wish you a speedy recovery!
\end{quote}

Neuter.

\begin{cousinweb}
\emph{\ru{выздоровление}} is built on \emph{\ru{выздороветь}} 'to recover', from \emph{\ru{здоровый}} 'healthy', with \emph{-\ru{ение}}; EtymDB has no entry for the noun.
\end{cousinweb}

\begin{grammarlens}[title={What you've built}]
One more word for telling, explaining and describing at B1.
\end{grammarlens}

\subsection*{Your turn}
\begin{itemize}
\raggedright
  \item \emph{Say it:} \emph{vyzdorovléniye}
  \item \emph{Say it:} \emph{vyzdorovléniye}, once more
  \item \emph{Say it:} a sentence of your own with \emph{vyzdorovléniye}
  \item \emph{Recall:} say the Russian for a surgeon, then the Russian for to recover, then say \emph{vyzdorovléniye} again
\end{itemize}

\begin{tcolorbox}[breakable,colback=teal!4,colframe=teal!35!black,title={Before you move on}]
What does \ru{выздоровление} mean? (\textquotedblleft{}Recovery\textquotedblright{}.) Say it once more.
\end{tcolorbox}
\section[\texorpdfstring{\ru{поправляться} (popravlyátsya)}{поправляться (popravlyátsya)}]{\ru{поправляться} (popravlyátsya) — to get better; to put on weight}

\label{lesson:RU-C490-popravlyatsya}



\begin{quote}
Before the new one: say the Russian for to recover, then the Russian for recovery.
\end{quote}

\subsection*{You'll want to know: \ru{поправляться}}
\textbf{\ru{поправляться}} (\emph{popravlyátsya}) — \textquotedblleft{}to get better; to put on weight\textquotedblright{}.

\begin{quote}
\ru{После} \ru{болезни} \ru{она} \ru{быстро} \ru{поправляется}.
\end{quote}

\begin{quote}
\emph{Pósle bolézni oná býstro popravlyáyetsya.}
\end{quote}

\begin{quote}
After her illness she is getting better quickly.
\end{quote}

Imperfective; perfective \ru{поправиться} (\ru{поправлюсь}). Also to put on weight.

\begin{cousinweb}
\emph{\ru{поправляться}} is the reflexive of \emph{\ru{поправлять}}, the partner of \emph{\ru{поправить}} 'to set right', built on \emph{\ru{править}} 'to put right' and \emph{\ru{правый}} 'right', which EtymDB traces to Proto-Slavic \emph{prav\ru{ъ}} 'right, correct'.
\end{cousinweb}

\begin{grammarlens}[title={What you've built}]
One more word for telling, explaining and describing at B1.
\end{grammarlens}

\subsection*{Your turn}
\begin{itemize}
\raggedright
  \item \emph{Say it:} \emph{popravlyátsya}
  \item \emph{Say it:} \emph{popravlyátsya}, once more
  \item \emph{Say it:} a sentence of your own with \emph{popravlyátsya}
  \item \emph{Recall:} say the Russian for to recover, then the Russian for recovery, then say \emph{popravlyátsya} again
\end{itemize}

\begin{tcolorbox}[breakable,colback=teal!4,colframe=teal!35!black,title={Before you move on}]
What does \ru{поправляться} mean? (\textquotedblleft{}To get better; to put on weight\textquotedblright{}.) Say it once more.
\end{tcolorbox}
\section[\texorpdfstring{\ru{глотать} (glotát)}{глотать (glotát)}]{\ru{глотать} (glotát) — to swallow}

\label{lesson:RU-C490-glotat}



\begin{quote}
Before the new one: say the Russian for recovery, then the Russian for to get better.
\end{quote}

\subsection*{You'll want to know: \ru{глотать}}
\textbf{\ru{глотать}} (\emph{glotát}) — \textquotedblleft{}to swallow\textquotedblright{}.

\begin{quote}
\ru{Мне} \ru{больно} \ru{глотать}.
\end{quote}

\begin{quote}
\emph{Mne bólno glotát.}
\end{quote}

\begin{quote}
It hurts to swallow.
\end{quote}

Imperfective; perfective \ru{проглотить} (\ru{проглочу}).

\begin{cousinweb}
EtymDB lists \emph{\ru{глотать}}, glossed 'to swallow', without a source; it is usually traced to Proto-Slavic \emph{gl\ru{ъ}tati}.
\end{cousinweb}

\begin{grammarlens}[title={What you've built}]
One more word for telling, explaining and describing at B1.
\end{grammarlens}

\subsection*{Your turn}
\begin{itemize}
\raggedright
  \item \emph{Say it:} \emph{glotát}
  \item \emph{Say it:} \emph{glotát}, once more
  \item \emph{Say it:} a sentence of your own with \emph{glotát}
  \item \emph{Recall:} say the Russian for recovery, then the Russian for to get better, then say \emph{glotát} again
\end{itemize}

\begin{tcolorbox}[breakable,colback=teal!4,colframe=teal!35!black,title={Before you move on}]
What does \ru{глотать} mean? (\textquotedblleft{}To swallow\textquotedblright{}.) Say it once more.
\end{tcolorbox}
\section[\texorpdfstring{\ru{сироп} (siróp)}{сироп (siróp)}]{\ru{сироп} (siróp) — syrup}

\label{lesson:RU-C490-sirop}



\begin{quote}
Before the new one: say the Russian for to get better, then the Russian for to swallow.
\end{quote}

\subsection*{You'll want to know: \ru{сироп}}
\textbf{\ru{сироп}} (\emph{siróp}) — \textquotedblleft{}syrup\textquotedblright{}.

\begin{quote}
\ru{Врач} \ru{выписал} \ru{сироп} \ru{от} \ru{кашля}.
\end{quote}

\begin{quote}
\emph{Vrach výpisal siróp ot káshlya.}
\end{quote}

\begin{quote}
The doctor prescribed cough syrup.
\end{quote}

Masculine.

\begin{cousinweb}
\emph{\ru{сироп}} was borrowed from French \emph{sirop}, which comes from Arabic \emph{šarāb} 'drink' (EtymDB).
\end{cousinweb}

\begin{grammarlens}[title={What you've built}]
One more word for telling, explaining and describing at B1.
\end{grammarlens}

\subsection*{Your turn}
\begin{itemize}
\raggedright
  \item \emph{Say it:} \emph{siróp}
  \item \emph{Say it:} \emph{siróp}, once more
  \item \emph{Say it:} a sentence of your own with \emph{siróp}
  \item \emph{Recall:} say the Russian for to get better, then the Russian for to swallow, then say \emph{siróp} again
\end{itemize}

\begin{tcolorbox}[breakable,colback=teal!4,colframe=teal!35!black,title={Before you move on}]
What does \ru{сироп} mean? (\textquotedblleft{}Syrup\textquotedblright{}.) Say it once more.
\end{tcolorbox}
